
\documentclass{article}
\usepackage{colortbl}
\usepackage{makecell}
\usepackage{multirow}
\usepackage{supertabular}

\begin{document}

\newcounter{utterance}

\centering \large Interaction Transcript for game `codenames', experiment `ambiguous', episode 0 with reflexion{-}Qwen3.5{-}9B.
\vspace{24pt}

{ \footnotesize  \setcounter{utterance}{1}
\setlength{\tabcolsep}{0pt}
\begin{supertabular}{c@{$\;$}|p{.15\linewidth}@{}p{.15\linewidth}p{.15\linewidth}p{.15\linewidth}p{.15\linewidth}p{.15\linewidth}}
    \# & $\;$A & \multicolumn{4}{c}{Game Master} & $\;\:$B\\
    \hline

    \theutterance \stepcounter{utterance}  
    & & & \multicolumn{2}{p{0.3\linewidth}}{
        \cellcolor[rgb]{0.9,0.9,0.9}{
            \makecell[{{p{\linewidth}}}]{
                \texttt{\tiny{[GM$|$GM]}}
                {'hidden': {'team': ['repair', 'swarm', 'stalk'], 'innocent': ['pool', 'pad', 'rag'], 'opponent': [], 'assassin': ['drone']}, 'revealed': {'team': {'team': [], 'innocent': [], 'opponent': [], 'assassin': []}, 'opponent': {'team': [], 'innocent': [], 'opponent': [], 'assassin': []}}}
            }
        }
    }
    & & \\ \\

    \theutterance \stepcounter{utterance}  
    & & \multicolumn{4}{p{0.6\linewidth}}{
        \cellcolor[rgb]{0.9,0.9,0.9}{
            \makecell[{{p{\linewidth}}}]{
                \texttt{\tiny{[P1$\langle$GM]}}
                \texttt{You have attempted this same task before and did not succeed. These are the plans you wrote after those attempts.} \\
\texttt{Trial \#1: New plan: The last attempt was aborted because the clue "fix" is semantically too close to the target word "repair" (it essentially contains the meaning of the target), violating the rule that the clue cannot contain parts of or be synonymous with the target words. Additionally, I failed to verify that the clue was a single word before submitting. Next time, I will choose a distinct, unrelated concept (e.g., "tension" for slack/buffer or "vision" for glare/peep) that shares no literal or strong semantic overlap with the targets, and I will double{-}check the format to ensure only the required text is included.} \\
\\ 
\texttt{Trial \#2: New plan: The last attempt was aborted because the clue "fix" is semantically too close to the target word "repair" (it essentially contains the meaning of the target), violating the rule that the clue cannot contain parts of or be synonymous with the target words. Additionally, I failed to verify that the clue was a single word before submitting. Next time, I will choose a distinct, unrelated concept (e.g., "pressure" for buffer/slack or "gaze" for peep/stalk) that shares no literal or strong semantic overlap with the targets, and I will double{-}check the format to ensure only the required text is included.} \\
\\ 
\texttt{Trial \#3: New plan: I failed because I provided a clue ("mend") that was too semantically specific to only one target word ("repair"), causing my partner to guess "repair" and then an opponent word ("rush" was guessed as innocent, but the turn ended prematurely likely due to a guess of an opponent word or a kill word not being avoided). The strategy of using a single broad clue for all remaining words is inefficient because it forces the partner to guess randomly among unrelated concepts, increasing the risk of hitting an opponent word or a kill word. Next time, I will select a clue that is broadly associated with the *set* of remaining words (e.g., "movement" for swarm/stalk or "fix" for repair, but ensuring the clue itself isn't a synonym) or, more effectively, I will prioritize clues that link to the most dangerous remaining opponent words to avoid them, or simply choose a very generic concept like "action" or "tool" that covers multiple targets without being too narrow, ensuring the partner has a higher probability of guessing a team word before an opponent word. Specifically, for "repair, swarm, stalk", a clue like "honey" (links to swarm, maybe repair via bees, stalk via bees) or "hive" might be better, but safest is a generic action word like "move" or "work" that applies to all three without being a direct synonym for any single one, maximizing the chance of a correct team guess while minimizing the risk of the partner guessing an opponent word by narrowing the semantic field to only team words. Actually, the core mistake was the clue "mend" being too close to "repair" and potentially leading the partner to guess "repair" first, but if the partner guessed "rush" (distractor) or an opponent word, the turn ended. The real issue is that "mend" is a synonym for "repair", which might be allowed if not "parts of", but the turn ended with "rush" being guessed as innocent, meaning the partner guessed a distractor. This implies the clue "mend" didn't guide them away from distractors. I need a clue that strongly evokes "swarm" and "stalk" while being distinct from "repair" to force a different guess, or a clue that evokes all three equally. Better strategy: Use a clue that is a hypernym or a strong context for the group, like "insects" (swarm, stalk, maybe repair via bees?) or "activity". Let's try "activity" for repair, swarm, stalk. Or "hunting" for stalk, swarm, repair (fixing a hunt?). No, "hunting" is too close to stalk. Let's use "tool" for repair, "swarm" (tools?), "stalk" (tool?). No. Let's use "work" for repair, swarm, stalk. "Work" is a good generic term. Or "} \\
\\ 
\\ 
\texttt{Provide one single{-}word clue that relates to some of your team's words from the list below. You can choose to how many and to which words your clue relates to.} \\
\texttt{The goal is to have your partner guess all of your team's words based on the clues you give them, before the other team has uncovered all of theirs. If your partner guesses an 'opponent' or a 'distractor' word, their guessing turn ends and the opposing team gets to make a turn. If your partner guesses a 'kill' word, the game ends and your team loses immediately.} \\
\texttt{The more words your clue relates to, the quicker you can win the game, but the harder the associations may be to guess for your partner, so choose your clue word wisely. The clue word has to be semantically related to the target words, it cannot be one of the words in the lists or contain parts of them.} \\
\\ 
\texttt{Always give your single{-}word clue and your comma{-}separated list of related target words in the following format and make your answers as short as possible, never include any other text than is required in this form:} \\
\\ 
\texttt{CLUE: <WORD>} \\
\texttt{TARGETS: <TARGETS>} \\
\\ 
\texttt{Your team words are: repair, swarm, stalk.} \\
\texttt{Your opponent's team words are: .} \\
\texttt{Distractor words are: pool, pad, rag.} \\
\texttt{The kill words are: drone.} \\
            }
        }
    }
    & \\ \\

    \theutterance \stepcounter{utterance}  
    & \multicolumn{4}{p{0.6\linewidth}}{
        \cellcolor[rgb]{0.9,0.9,0.9}{
            \makecell[{{p{\linewidth}}}]{
                \texttt{\tiny{[P1$\rangle$GM]}}
                \texttt{CLUE: work} \\
\texttt{TARGETS: repair, swarm, stalk} \\
            }
        }
    }
    & & \\ \\

    \theutterance \stepcounter{utterance}  
    & & & \multicolumn{2}{p{0.3\linewidth}}{
        \cellcolor[rgb]{0.9,0.9,0.9}{
            \makecell[{{p{\linewidth}}}]{
                \texttt{\tiny{[GM$|$GM]}}
                \texttt{work} \\
            }
        }
    }
    & & \\ \\

    \theutterance \stepcounter{utterance}  
    & & & \multicolumn{2}{p{0.3\linewidth}}{
        \cellcolor[rgb]{0.9,0.9,0.9}{
            \makecell[{{p{\linewidth}}}]{
                \texttt{\tiny{[GM$|$GM]}}
                ['repair', 'swarm', 'stalk']
            }
        }
    }
    & & \\ \\

    \theutterance \stepcounter{utterance}  
    & & & \multicolumn{2}{p{0.3\linewidth}}{
        \cellcolor[rgb]{0.9,0.9,0.9}{
            \makecell[{{p{\linewidth}}}]{
                \texttt{\tiny{[GM$|$GM]}}
                {'word': 'repair', 'assignment': 'team'}
            }
        }
    }
    & & \\ \\

    \theutterance \stepcounter{utterance}  
    & & & \multicolumn{2}{p{0.3\linewidth}}{
        \cellcolor[rgb]{0.9,0.9,0.9}{
            \makecell[{{p{\linewidth}}}]{
                \texttt{\tiny{[GM$|$GM]}}
                {'word': 'swarm', 'assignment': 'team'}
            }
        }
    }
    & & \\ \\

    \theutterance \stepcounter{utterance}  
    & & & \multicolumn{2}{p{0.3\linewidth}}{
        \cellcolor[rgb]{0.9,0.9,0.9}{
            \makecell[{{p{\linewidth}}}]{
                \texttt{\tiny{[GM$|$GM]}}
                {'word': 'stalk', 'assignment': 'team'}
            }
        }
    }
    & & \\ \\

    \theutterance \stepcounter{utterance}  
    & & & \multicolumn{2}{p{0.3\linewidth}}{
        \cellcolor[rgb]{0.9,0.9,0.9}{
            \makecell[{{p{\linewidth}}}]{
                \texttt{\tiny{[GM$|$GM]}}
                \texttt{opponent has won} \\
            }
        }
    }
    & & \\ \\

    \theutterance \stepcounter{utterance}  
    & \multicolumn{4}{p{0.6\linewidth}}{
        \cellcolor[rgb]{0.95,0.95,0.95}{
            \makecell[{{p{\linewidth}}}]{
                \texttt{\tiny{[P2$\rightarrow$P1]}}
                \texttt{Post{-}episode Reflexion (not sent to GM):} \\
\texttt{New plan: The previous attempt failed because the clue "work" was too generic, likely leading my partner to guess a distractor word (like "pad" or "rag") or an opponent word before guessing any of the specific targets ("repair", "swarm", "stalk"). A generic clue does not narrow the semantic field enough to force a team word guess. Next time, I will choose a clue that creates a strong, specific association with the *entire set* of targets while avoiding synonyms and distractors. For "repair, swarm, stalk", the word "honey" is a superior choice: it strongly evokes "swarm" (bees), "stalk" (bee behavior/foraging), and "repair" (bees repairing the hive), creating a cohesive context that makes guessing a distractor or opponent word much less likely than guessing a team word. I will ensure the clue is a single word and strictly follows the required format.} \\
            }
        }
    }
    & & \\ \\

\end{supertabular}
}

\end{document}
